\oldpage[IV-2]{195}

\begin{env}[7.3.10]
\label{IV.2.s7.3.10-fr}
Diremos que una propiedad $\boldsymbol P(Z,k)$ es del \emph{primer
tipo} si la propiedad $\boldsymbol Q(A,k)$ que sirve para definir
$\boldsymbol P$ es una propiedad $\boldsymbol R(A)$ \emph{en la que no
interviene $k$} y se cumple cuando $A$ es un \emph{cuerpo}; diremos que
$\boldsymbol P(Z,k)$ es del \emph{segundo tipo} si la propiedad
$\boldsymbol Q(A,k)$ es de la forma
\begin{quote}
«Para toda extensión $k'$ de tipo finito de $k$ y todo punto $z'$ de
$\operatorname{Spec}(A\otimes_k k')$ situado sobre el punto cerrado de
$\operatorname{Spec}(A)$, se cumple $\boldsymbol R(\mathscr O_{z'})$».
\end{quote}
Se sigue suponiendo que $\boldsymbol R(A)$ se cumple cuando $A$ es un
cuerpo. Es claro que, en los ejemplos de (7.3.8), las propiedades (i)
a (vii) son del primer tipo, y las propiedades (iv$'$) a (vii$'$), del
segundo tipo.

Si se retoman los razonamientos de (6.6.1) y (6.8.3), se ve que, cuando
$\boldsymbol P$ es del primer o del segundo tipo para una propiedad
$\boldsymbol R(A)$, \emph{las condiciones $(\mathrm P_{\mathrm I})$ y
$(\mathrm P_{\mathrm{II}})$ se deducen de las siguientes condiciones
sobre $\boldsymbol R$}:

$(\mathrm R_{\mathrm I})$ Sean $X$ e $Y$ dos preesquemas localmente
noetherianos y $f:X\to Y$ un morfismo \emph{regular}; entonces, para
todo $x\in X$, la propiedad $\boldsymbol R(\mathscr O_{f(x)})$ implica
la propiedad $\boldsymbol R(\mathscr O_x)$.

$(\mathrm R_{\mathrm{II}})$ Sea $\varphi:A\to B$ un homomorfismo local
de anillos locales noetherianos que haga de $B$ un $A$-módulo \emph{plano};
entonces $\boldsymbol R(B)$ implica $\boldsymbol R(A)$.

Además, la propiedad $(\mathrm P_{\mathrm{III}})$ se deduce de que
$\boldsymbol R(k)$ se cumple por hipótesis para todo cuerpo $k$; por
último, cuando $\boldsymbol P$ es del segundo tipo, la condición
$(\mathrm P_{\mathrm{IV}})$ se deduce de la definición de $\boldsymbol P$
y de la transitividad de las extensiones de tipo finito.
\end{env}
